% 本文件由 LaTeX 排版 Agent 根据有效 translation.json 自动生成。
% author=Tertullian; work_id=0408; title=论斋戒

\chapter{论斋戒}

% structure: p0001 source=h1 target=omitted reason=page-title-already-rendered-by-parent

\Emph{反对属血气者。}\par

% structure: p0003 source=h3 target=section reason=relative-heading-rank; base=h3

\section{第一章 贪食与情欲的关联。属血气者反对孟他努派的理由}

我若见属血气者只沉溺于纵欲——这诱惑他们多次嫁娶——而不同时因贪食而膨胀，以致憎恶斋戒，我必诧异。因为情欲无贪食，无疑可视为怪异现象；此二者如此相连而结合，以致若有可能将其分开，私处也不会依附于肚腹本身，而应在他处。试观身体：这些器官的区域本为一体。简言之，恶习的次序与肢体的排列相称。先是肚腹；随后一切其他纵欲的材料立即从属于美味：通过贪食之爱，不洁之爱得以通行。因此，我认出那\Emph{属血气的}信德——它完全由肉身构成，所以它既贪多食，亦贪多婚——它理应指控那\Emph{属神的}纪律，因为后者按己力对抗它，同样在这方面的节制中：即在食欲上加以缰绳，或有时不进食，或迟进食，或干进食，正如在情欲上加以缰绳一样，只允一夫一妻制。\par

与这样的人争辩实在烦人：争论那些连辩护本身都冒犯谦逊的话题，实在羞愧。因为我若不指责他们的对手，又怎能维护贞洁与节制呢？那些对手，我姑且一言以蔽之：就是属血气者外在和内在的\OriginalTerm{香肠}{botuli}。正是这些与护慰者挑起争论；正是为此，新预言被拒绝：并非因为孟他努、普黎斯奇拉和玛克西弥拉宣讲另一位天主，也非他们将耶稣基督与天主分开，亦非他们推翻任何信德或望德的规则，而是因为他们明确教导更频繁的斋戒而不是嫁娶。关于嫁娶的限度，我们已经出版了一篇为一夫一妻制辩护的文章。现今我们的争战是关于次要的（或毋宁说首要的）节制，即在饮食的克制方面。他们指控我们遵守自己的斋戒；通常将我们的站礼延长至傍晚；也遵守干食，使食物不受任何肉、任何汁液及任何特别多汁果实的湿润；也不吃或饮用任何带酒味的东西；还有与干食相配的沐浴戒除。因此他们不断以新奇指责我们；关于其非法性，他们立下一条惯例，即必须判定为\Emph{异端}（若争论之点是人的臆测），或宣布为\Emph{假预言}（若是属神的宣告）；无论如何，我们这些主张回归者都听到绝罚的判决。\par

% structure: p0006 source=h3 target=section reason=relative-heading-rank; base=h3

\section{第二章 属血气者从法律、福音、宗徒大事录、书信及异教习俗引用的论据}

因为，就斋戒而言，他们以天主所指定的确定日子来反驳我们：例如，在\BibleBook{}{肋未纪}中，主命令\Person{梅瑟}将七月第十日定为赎罪日，说：\BibleQuote{这日应为你们至圣之日，你们要克苦己心；任何在那日不克苦己心的灵魂，应从其人民中铲除。}无论如何，他们认为在福音中，那些\BibleQuote{新郎被带走}的日子已被确定为斋戒的日子；而现在这些是基督徒斋戒唯一合法的日子，法律和先知的古制已被废除：因为无论何处合乎他们的意愿，他们承认\BibleQuote{法律及先知到\Person{若翰}为止}的含义。因此，他们认为，关于未来，斋戒应由新纪律随意遵守，出于选择而非命令，根据各人的时间和需要：而且这曾是宗徒们的遵守方式，他们没有给普遍遵守的确定斋戒加上其他轭，也没有同样给\OriginalTerm{站斋}{Stations}加上轭，这些站斋固然有它们自己的日子（第四日和第六日），但仍按个人判断广泛施行，既不受既定诫命的法律约束，也不延长至一天的最后一小时，因为甚至祈祷也通常以第九时辰结束，遵照\BibleBook{}{宗徒大事录}所记载的\Person{伯多禄}的榜样。然而，他们认为\OriginalTerm{干食}{Xerophagies}是一种刻意修行的新名称，非常接近异教迷信，如同以某些食物的限制来净化\OriginalTerm{阿匹斯}{Apis}、\OriginalTerm{伊西斯}{Isis}和\OriginalTerm{大母神}{Magna Mater}的禁欲严规；而信仰在基督内是自由的，即使对犹太法律也不欠特定肉食的节制，这已被宗徒一次而永远地允许进入整个肉市场——（宗徒），他憎恶那些像禁止婚姻一样，命令我们戒绝天主所造食物的人。因此，他们认为我们甚至在那时已被预告为\BibleQuote{在末期要离弃信德，听从那迷惑世界的邪神，良心被烙铁烙过，随从说谎者的道理。}被烙？请问，被什么火？我想，是那火引导我们一再缔结婚姻和每日烹饪正餐！同样，他们断言我们与迦拉达人一同分享（宗徒）尖锐的责备，作为\BibleQuote{遵守日子、月份和年份的人。}同时，他们当面嘲笑我们，因为\Person{依撒意亚}也曾权威地宣告：\BibleQuote{这样的斋戒不是主所拣选的}，即不是禁食，而是正义的行为，他在那里附加了这些；并且主自己在福音中对各类关于食物的挑剔给出了简洁的回答：\BibleQuote{不是进入口内的使人污秽，而是从口里出来的才使人污秽}；同时他自己也曾吃喝以至于如此被人注意：\BibleQuote{看，一个贪吃嗜酒的人}；（最后），宗徒也教导说\BibleQuote{食物不能使我们得到天主的赞许；因为我们吃，也不多；不吃，也不少}。\par

借这些和类似段落的工具，他们最终诡辩到如此程度，以至于每个比较贪恋口腹的人都能认为禁食、减食或推迟进食的义务是多余的、不那么必要的，因为\Quote{天主}毕竟\Quote{喜欢正义和纯朴的行为}。我们知道肉体便利的劝勉言词的性质，很容易说：\Quote{我必须全心相信；我必须爱天主，并爱人如己：因为\InnerQuote{全部法律和先知都系于这两条诫命}，而不是系于我肺腑的空虚。}\par

% structure: p0009 source=h3 target=section reason=relative-heading-rank; base=h3

\section{第三章 斋戒原则追溯至其最早源头}

因此，在进一步讨论之前，我们必须确认这一（原则），它有被暗中颠覆的危险；即，你所谈论的这种\Quote{空虚}在天主眼中究竟有何价值：首先，以此方式赚取天主恩宠的理据本身是从何而来的。因为当从起初就可追溯的理据的权威显现出来时，遵守的必要性就会被承认。\par

亚当从天主领受了不可品尝\\BibleQuote{分辨善恶之树}的律法，并附有品尝后死亡的判决。然而，就连他（亚当）自己，当时在神魂超拔中曾预言性地解释过那涉及基督与教会的\\BibleQuote{伟大奥迹}之后，退回到\\OriginalTerm{属魂之人}{Psychic}的状态，不再能\\BibleQuote{领会属神之事}，便更顺从肚腹而非天主，更留意肉食而非诫命，为口腹之欲出卖了救恩！总之，他吃了，便灭亡了；他本可得救，如果他宁愿从一棵小树上禁食；因此，从这早期起，\\Emph{\\OriginalTerm{属血肉的信德}{animal faith}}就能认出自己的种子，从此推断出它对肉欲的嗜好和对属灵事物的拒绝。因此，我坚持认为，从一开始那杀人的喉咙就应受到饥饿的折磨和惩罚。即使天主没有命令任何禁食的规诫，但通过指出亚当被杀的原因，那位显示了过犯的天主已将过犯的补救方法留给了我的理智。即使没有命令，我也会在可能的方式和时间里，习惯性地把食物视为毒药，并服用解药——饥饿；藉此清除死亡的原始原因——这原因与我的出生同时传给了我；确信天主愿意那相反之事祂所不愿的，并足够相信，克己的操劳会中悦于祂，因我从祂那里应已明白，\\Emph{不}节制之罪已被定罪。此外：既然祂自己既命令禁食，又呼唤\\BibleQuote{破碎的灵魂}——当然，是由于饮食的紧迫——为\\BibleQuote{祭献}；谁还会怀疑所有饮食克苦的理据是：藉由重新禁止食物和遵守诫命，原始的罪如今可得赎罪，好使人通过自己得罪天主的同一原因材料向天主做出赔补，即通过禁止食物；并且，仿效地，饥饿可能重新点燃救恩，正如饱足曾熄灭它，为了一个不法的而鄙弃更合法的（满足）？\par

% structure: p0012 source=h3 target=section reason=relative-heading-rank; base=h3

\section{第四章 提出的异议：为何洪水后合法食物的限制被扩大了？回答}

这理据始终在天主的眷顾眼中——祂调节一切以适应时代的需要——免得任何反对者为了破坏我们的主张而说：\\Quote{为何那时天主没有立即制定某种明确的饮食限制？不，为何祂反而扩大了许可？因为起初，祂只将草本和树木的食物赐给人：\\BibleQuote{看，我赐给你们地上一切结种子的草，和一切树上所结有核的果子，都作为你们的食物。}然而，后来在向诺厄列举\\BibleQuote{地上一切野兽、天上一切飞鸟、地上一切爬行物、海里的鱼，以及一切爬虫}的隶属之后，祂说：\\BibleQuote{它们都作为你们的食物：如同青菜一样，我将它们全部赐给你们；但带血的肉，即带生命之血的肉，你们不可吃。}因为正是藉着这一点，祂只禁止吃那生命尚未通过血倾流的肉，显然祂已准许吃所有其他肉。}对此我们答辩：这不适合使人负担任何进一步的特殊禁食律法，因为人刚刚才显明不能忍受如此轻的禁令——即单单一枚果实；因此，放松了缰绳之后，他要藉自己的自由得以坚强；同样，在洪水之后，在人类的重新塑造中，（如同之前），\\Emph{一条}规律——禁血——就已足够，而准许使用所有其他东西。因为主已经通过洪水显示了他的审判；而且，还通过\\BibleQuote{向兄弟的手、向一切野兽的手索讨血}发出了威吓的警告。如此，预先施行审判的公义，祂发出了自由的原料；通过允许准备纪律的萌芽；允许一切，以便取走一些；意思是要\\BibleQuote{更多索讨}，如果祂曾\\BibleQuote{更多交托}；命令禁食，因为祂预见了放纵：为的是（如我们所说）原始的罪能藉着在更大自由的机会中施行更大禁食而更得赎罪。\par

% structure: p0014 source=h3 target=section reason=relative-heading-rank; base=h3

\section{第五章 转向以色列人历史，\\Person{戴尔都良}表明贪食如亚当之罪同样显著。因此设置了肋未法律的限制。}

最后，当一亲近的百姓开始被天主特选归己，而人的恢复也得以尝试时，便颁布了一切法律与纪律，包括限制食物的规定；某些物被禁为不洁，好使人通过在特定事上持续禁食，最终更易忍受彻底的斋戒。因为起初的百姓也重现了最初之人的罪行，被发现有更倾向肚腹而非天主的倾向——当他们被天主\BibleQuote{以大能的手和崇高的臂}从埃及奴役的严酷中拔出后，竟被看见成为那地的主人，被预定进入\BibleQuote{流奶与蜜之地}；但立刻，因四周荒凉景象而绊倒，叹息着失去的埃及饱足之乐，他们向梅瑟和亚郎抱怨说：\BibleQuote{但愿我们被主击中心脏，死在埃及地，那时我们常坐在肉锅旁，吃饼吃到饱足！你们为何领我们出到这些旷野，要叫这会众饿死呢？}由于同样的肚腹偏好，他们最终注定要哀叹自己同样的铅——和天主（能力）的目击者——他们因对肉的贪恋和对埃及丰足的回忆，不断激怒他们：\BibleQuote{谁给我们肉吃呢？我们想起了在埃及常可随意吃的鱼，还有黄瓜、西瓜、韭菜、葱和蒜。但现在我们的灵魂枯干：除了玛纳，我们眼中别无他物！}他们就这样（像属血肉者一样）也觉得天使的干食之饼不合口味：他们更喜欢蒜和葱的香气胜过天堂的芬芳。因此，从如此忘恩的人那里，一切更惬意和开胃之物都被撤去，为要同时惩罚贪食和操练节制，使前者受谴责，后者被实际学到。\par

% structure: p0016 source=h3 target=section reason=relative-heading-rank; base=h3

\section{第六章　斋戒与进食的肉体倾向考——梅瑟和厄里亚的实例}

如今，若我们追溯至原初经验以说明天主为何规定、我们为何（为天主之故）当规定饮食限制，这未免有些冒昧；那就让我们求问共同的良心吧。自然本身会清楚告诉我们，当她在\Emph{之前}进食和饮酒、我们的唾液尚处于纯真状态时，以处理神圣之事的感官，使我们具备何等能力参与事务——（是否）以更警觉的头脑、更活泼的心，远胜于当我们内那整个寓所塞满食物、浸满酒液、发酵以作排泄之预备时，已变成厕所的预谋之所，在那里，显然没有什么比随后立即发生的情欲滋味更临近的了。\BibleQuote{百姓吃喝，起来玩耍。}请理解圣经含蓄的语言：\Quote{玩耍，}若非不端庄，就不会被谴责。另一方面，有多少人当记忆的座位被占据、智慧的肢体受阻时仍不忘虔敬呢？在那时——即人惯于忘记自己之时——没有人能恰当、适宜、有益地记念天主。一切纪律之食物要么杀死，要么伤害。若主自己在责备以色列人忘记时，不将其原因归于\Quote{饱足}，我便是说谎者：\BibleQuote{（我的）所爱者变得肥厚、肥胖、鼓胀，完全离弃了造他的天主，离开了主他的救主。}简言之，在同一部\WorkTitle{申命纪}中，祂在叮嘱防范同一原因时说：\BibleQuote{恐怕你吃了——喝了，建造了美好的房屋，羊群牛群增多，（你的）金银增多——你的心就高傲，忘记主你的天主。}祂将贪食的过度置于财富的败坏力量之前，而财富本身正是贪食的促成者。通过它们，的确，\BibleQuote{百姓的心变得迟钝，免得他们用眼看见，用耳听见，用心明白，}这个心被\Quote{脂油}所阻塞——天主明确禁止食用那些，教导人不要专注肚腹。\par

另一方面，那人的\Quote{心}习惯性地被找到是\Quote{高举}而非养肥的，他在四十昼夜中维持了超越人性的守斋，同时属神的信德将力量提供（给身体），他亲眼看见天主的荣耀，亲耳听见天主的声音，亲心领悟天主的法律；那时祂甚至（借着经验）教导他：人生活不只靠饼，也靠天主口中所发的一切言语；因为那子民虽比他肥胖，却不能恒常注视甚至梅瑟本人，因他是以天主为食的，也不能注视他的瘦弱，因它已因天主的荣耀而饱足！因此，配得地，即使还在肉身内，上主就向祂自己显示给他，他成了祂自己守斋的同工，正如对厄里亚一样。因为厄里亚也由此事首先已将自己充分奉献于守斋：他咒诅了饥荒。\BibleQuote{上主永生活着，}他说，\BibleQuote{在我站在祂面前时，若在这几年中有露水，有降雨。}随后，因躲避威胁他的\Person{依则贝尔}，在天使唤醒他时仅吃了一餐食物和饮料后，他也在四十昼夜的空间里，肚腹空空，口干舌燥，到达了\Place{曷勒布山}；在那里，他使一个洞穴成为他的客栈，他是以何等熟悉的与天主的会晤被接待！\BibleQuote{厄里亚，你（在）这里做什么？}这声音远比\BibleQuote{亚当，你在哪里？}更友善，因为后者的声音是对一个吃饱的人发出威胁，前者则安慰一个守斋的人。这就是限制的饮食的殊荣：它使天主与人成为帐幕同伴——确实是同伴与同伴！因为如果永恒的天主不会饥饿，正如祂通过\Person{依撒意亚}所证实的，那么这就是人变得与天主同等的时候，即当他不需要食物而生活时。\par

% structure: p0019 source=h3 target=section reason=relative-heading-rank; base=h3

\section{第七章 旧约中有利于守斋的更多例子}

因此我们已经继续举例，为了借其有益的效能，揭示这项职责的力量，它能平息天主的愤怒，使祂与人和好。\par

以色列在通过\Person{撒慕尔}聚集于\Place{米兹帕}取水之前犯了罪；但他们如此迅速地以守斋洗刷了罪，以致战事的危险同时（随水泼地）被他们驱散。正当\Person{撒慕尔}献全燔祭时（我们绝不知道天主的仁慈比百姓的节制更能获得什么），外敌正前进作战，就在那里\BibleQuote{上主以巨雷向敌人怒吼，他们混乱了，在以色列面前成堆倒下；以色列人从米兹帕出来，追赶敌人，击杀他们直到贝特曷龙}——未进食的（追赶）进食的，未武装的追赶武装的。这就是那些\Quote{为天主守斋}者的力量。因为，为此上天作战。你们拥有一个条件，据此将赐予神圣的护佑，甚至对属灵的战争也是必要的。\par

同样，当亚述王\Person{散乃黑黎布}已攻取数城之后，通过\Person{辣布沙克}向以色列射出亵渎和威吓时，没有别的（除了守斋）使他改变主意，把他打发到雇土去。此后，还有什么比\Person{希则克雅}王的谦卑更能借天使之手消灭他军队中的十八万四千人呢？如果这是真的（确实是），他一听到关于敌人凶暴的宣告，就撕裂衣服，披上麻衣，并吩咐长老司祭们同样穿着，通过\Person{依撒意亚}亲近天主——守斋当然是他祈祷的随行陪伴。因为危急没有时间进食，麻衣也不关心饱足的精致。饥饿总是哀悼的伴侣，正如喜乐是饱足的附属。\par

借着这哀悼的伴侣和（这）饥饿，甚至那有罪的城\Place{尼尼微}也从预言的毁灭中被释放。因为悔罪已充分促进了守斋，在三天空间内持续，甚至使那些未激怒天主的牲畜也绝食。\Place{索多玛}和\Place{哈摩辣}如果守了斋，也会逃脱。这救济甚至被\Person{阿哈布}承认。当他在违命和偶像崇拜，以及\Person{纳波特}被\Person{依则贝尔}因葡萄园而杀害之后，\Person{厄里亚}斥责他说：\BibleQuote{你杀了人，又占有了产业吗？在狗舔了纳波特血的地方，你的血也必被舔。}——他\BibleQuote{谦卑自己，将麻衣披在身上，守斋，并睡在麻衣中。然后上主的话传给厄里亚：你看见阿哈布在我面前战栗吗？因为他战栗，我不在他有生之年降祸，但在他儿子的日子，我必降祸}——（他的儿子），他没有守斋。因此，向天主守斋是敬畏的工作；借助它，\Person{厄耳卡纳}的妻子\Person{亚纳}，先前不孕，求告时也容易从天主获得了她空腹（没有食物）的腹中充满一个儿子，甚至一位先知。\par

禁食不仅仅带来本性的改变、危险的避免、罪过的消除，同样还能使禁食者从天主那里获得 \OriginalTerm{认识奥秘}{recognition of mysteries} 的能力。请看\Person{达内尔}的例子。关于巴比伦王的梦，所有的智者都感到困惑；他们断言，若无外援，人类技巧无法参透。唯独\Person{达内尔}倚靠天主，知道如何配得天主的垂青，他要求三天的期限，与他的弟兄一同禁食；他的祈祷如此受赞赏，他便通晓了梦的次序和意义；暴君的智者得以保全性命；天主受到光荣；\Person{达内尔}得到尊荣。他命中注定后来也得到天主同样大的恩宠：在\Person{达理阿}王元年，当他仔细反复思索\Person{耶肋米亚}所预言的时期后，他面朝天主，禁食、披麻蒙灰。随即，天使被派到他那里，立刻声明这就是天主赞许的原因：他说：\BibleQuote{我来了，为向你说明，因为你是个可怜人}——即因禁食而可怜。如果他在天主面前是\BibleQuote{可怜的}，那么他在狮子坑中对狮子来说就是可畏的；在那里，他禁食六天，有天使为他预备早餐。\par

% structure: p0025 source=h3 target=section reason=relative-heading-rank; base=h3

\section{第八章 来自新约的同类例子}

我们也提出我们其余的（证据）。因为我们急于转向近代的证明。在福音的门槛上，女先知\Person{亚纳}，\Person{法奴耳}的女儿，\BibleQuote{她既认出了婴孩主，便向那些期待以色列得救赎的人宣讲了许多关于祂的事}，在长久专一守寡的卓越殊荣之后，还额外得到\BibleQuote{禁食}的见证；以此指出，教会侍奉者应有的职责，以及（也指出）基督被那些一婚且常禁食的人所理解，无人能及。\par

此后，主亲自以自己的禁食祝圣了自己的圣洗（并在他自己的圣洗中，祝圣了所有人的圣洗）；他有能力使\BibleQuote{石头变成饼}，或者说，使\Place{约但河}流出酒来，假如他是那样一个\BibleQuote{贪食好酒的人}。反倒，借着轻视食物的德行，他正将\BibleQuote{新人}引入对\BibleQuote{旧人}的\BibleQuote{严厉处理}，为要向魔鬼显明那新人（魔鬼又借着\Emph{食物}试探他）胜过饥饿的全部力量。\par

此后，他为禁食制定了一条法律——它们应当\BibleQuote{没有忧愁}：因为有益的事为何要忧愁呢？他也教导说，禁食是与更凶恶的恶魔争战的武器：因为若同一种操作是不义之神的灵进入的通道，那么它成为圣神进入的通道又有何奇怪呢？最后，既然承认在百夫长\Person{科而乃略}身上，甚至在\Emph{圣洗之前}，圣神可敬的恩赐连同预言的恩赐已急速降临，我们看到，\Emph{他的禁食}已蒙垂听，我认为，此外，宗徒在\WorkTitle{格林多后书}中，在他的劳苦、危险和艰难之中，在\BibleQuote{又饥又渴}之后，也列举了\BibleQuote{禁食}，\BibleQuote{多次}。\par

% structure: p0029 source=h3 target=section reason=relative-heading-rank; base=h3

\section{第九章 从完全禁食到部分禁食和\OriginalTerm{干食}{xerophagiae}}

这种饮食限制的主要类别已经可以预先判断次等克己行为也同样依其程度有用或必要。因为排除某些种类的食物就是部分禁食。因此，让我们审视\OriginalTerm{干食}{xerophagiae}的新颖性或虚浮性，看看我们是否在其中也发现一种既古老又有效的宗教行为。我回到\Person{达内尔}和他的同伴，他们偏爱蔬菜和水的饮食，胜过王室的菜肴和酒器，并因此被发现\BibleQuote{更俊美}（免得有人还要为他瘦弱的身体担忧！），而且还在精神上受到培养。因为天主赐给这些青年各种文学的知识和智慧，赐给\Person{达内尔}各种言语、异梦和智慧；这智慧要使他在这件事上也聪明——即如何从天主获得\OriginalTerm{认识奥秘}{recognition of mysteries}的方法。最后，在波斯王\Person{居鲁士}第三年，当他陷入对异象的仔细反复默想时，他采取了另一种谦卑的形式。他说：\BibleQuote{那些日子，我达内尔哀悼了三周：我没有吃美味的面包，肉和酒没有进入我的口，我没有用油抹身，直到三周满了。}满了之后，一位天使（从天主）被派来，如此对他说：\BibleQuote{达内尔，你是个可怜人；不要怕：因为从你第一天起，你就将自己的灵魂献给回忆和在天主面前谦卑，你的话已蒙垂听，我因你的话而来。}于是，\BibleQuote{可怜}的景象和干食的谦卑驱走恐惧，吸引天主的耳朵，使人成为奥秘的主人。\par

我也回到厄里亚。当乌鸦素常以\Quote{面包和\Emph{肉}}供给他时，为何后来在犹大的贝尔舍巴，那位天使唤醒他之后，无疑只提供给他\Emph{单独}面包和水呢？难道是缺乏乌鸦，不能更丰盛地供给他吗？或者对那位\Quote{天使}来说，从国王宴席的某个平锅中，带走一位侍从及其丰盛的托盘，转移给厄里亚，像收割者的早餐被送入狮子洞、呈现给饥饿的达尼尔那样，是困难的吗？但是必须树立一个榜样，教导我们在压力、迫害和任何困难之时，必须靠\OriginalTerm{干食主义}{xerophagies}生活。达味就是用这样的食物表达了他的\OriginalTerm{认罪}{exomologesis}；\BibleQuote{吃灰尘如同吃饼}，即干如灰尘的面包：\BibleQuote{又将他的饮料与眼泪混合}——当然，代替酒。因为\Emph{戒酒}本身有荣耀的标志：这种戒酒曾将撒慕尔奉献给天主，并将亚郎祝圣给天主。因为关于撒慕尔，他的母亲说：\BibleQuote{清酒和烈酒他不应喝}，因为她在向天主祈祷时也是这样的状况。主对亚郎说：\BibleQuote{清酒和烈酒你不可喝，你和你的儿子在你之后，几时你们进入会幕，或走近祭坛，免得死亡}。确实如此，凡在教会中服务却不清醒的人，将\BibleQuote{死}。同样，在近代，祂责备以色列：\BibleQuote{你们给我的圣者酒喝}。而且，这种对饮料的限制是干食的一部分。无论如何，凡是天主命令或人发誓戒酒的地方，也应当理解为对食物的限制，作为对饮料的一种正式规范。因为饮料的性质与食物的性质相对应。一个人不太可能向天主奉献\Emph{一半}的食欲；在水中节制，在肉类中放纵。此外，宗徒是否了解干食主义？——他多次实践更大的刻苦：\BibleQuote{饥饿、口渴、多次禁食}，他曾禁止\BibleQuote{醉酒和宴乐}——我们甚至从他弟子弟茂德的例子有充分证据；当他劝诫弟茂德\BibleQuote{为了他的胃和常有的软弱}而用\BibleQuote{一点酒}时，弟茂德戒酒并非出于规则，而是出于虔诚——否则这习惯本会有益于他的胃——正是这一点，他建议戒酒是\BibleQuote{相称天主的}，而他出于\Emph{必要性}的理由，劝阻了戒酒。\par

% structure: p0032 source=h3 target=section reason=relative-heading-rank; base=h3

\section{第十章 论\OriginalTerm{斋期}{Stations}与祈祷的时辰}

同样地，他们以新奇为由批评我们的\OriginalTerm{斋期}{Stations}是\Emph{被规定的}；有些人还批评它们习惯性地拖延得太晚，说这一义务也应当自由选择遵守，不应该持续到第九时辰以后——当然，他们是从自己的实践得出规则。好了，关于\Emph{命令}的问题，我将一劳永逸地给出适合所有情况的答复。现在，回到这个特定情况的问题——即时间限制的问题——我必须首先问他们，他们从哪里得到这条在第九时辰结束斋期的规定？如果是因为我们读到伯多禄和他同在的人进入圣殿\BibleQuote{在第九时辰，祈祷的时辰}，谁能向我证明他们那天守了斋期，从而把第九时辰解释为斋期结束和解除的时间？相反，你更容易发现伯多禄在\Emph{第六}时辰，为了进食，先上了屋顶祈祷；所以一天的\Emph{第六}时辰也许更适合作为这一义务的界限，在伯多禄的例子中，祈祷后似乎就结束了这一义务。此外，既然在路加同一部记述中，\Emph{第三}时辰被证明是祈祷的时辰，在那个时辰，那些领受了圣神初始恩赐的人被认为是醉汉；\Emph{第六}时辰，伯多禄上了屋顶；\Emph{第九}时辰，他们进了圣殿：我们为什么不明白，我们绝对必须随时随地、无时无刻地祈祷，然而这三个时辰，在人类事务中更为显著——它们划分一天，分别各项事务，在公众耳中回响——在神圣的祈祷中也一直具有特殊庄严性？这种信念也得到了达尼尔每日祈祷三次这一证实性事实的认可；当然，通过某些指定时辰的例外，没有别的，正是更显著且后来成为宗徒传统的时辰——第三、第六、第九。因此，据此我将断言，伯多禄也是被古代习俗引导去遵守第九时辰的，在第三个特定间隔——最后的祈祷间隔——祈祷。\par

此外，我们提出这些（论点）是为了那些认为自己是依照伯多禄的榜样行事、却对其一无所知的人；并非我们轻视第九时辰——在每周的第四日和第六日，我们最隆重地尊崇它——而是因为，对于那些基于传统而遵守的事，我们负有责任提出更为相称的理由，因为它们缺乏圣经的权威，直到藉着某种显著的天上恩赐，它们或被确认或被纠正。宗徒说：\Quote{若有什么你们不明白的事，主会启示给你们。}因此，暂且不论这一切事的确认者——护慰者，即普世真理的引导者——请查考我们中间是否有一个更为相称的理由来遵守第九时辰；这样，如果伯多禄在那时遵守了站岗，我们的这个理由也必须归给他。因为（这实践）源自主的死亡；虽然这死亡应当时刻纪念，不分时辰，但我们在那时更深刻地受到提醒去纪念它，根据\Quote{站岗}这一名称的实际意义。因为即使士兵从不忘记他们的军誓，也对站岗更加敬重。因此，必须将那\Quote{压力}保持到那一时辰，即从第六时辰起，整个世界陷入普遍的黑暗，为它死去的主举行了悲哀的职守；这样，当宇宙重获阳光时，我们也可以回到喜乐中。如果这一点更符合基督宗教的精神，因为它更光荣地庆祝基督的光荣，我同样能够从同一事件顺序中，确定\Quote{站岗的晚延}的条件：即我们要禁食到很晚，等候主埋葬的时刻，那时若瑟取下了他所求的身体并安葬了他。由此（推知），仆人连在主人之前进食，甚至是不敬神的。\par

但是，到此为止，让我们满足于在\Quote{论证性的}挑战上交锋；正如我所做的，用猜测反驳猜测，然而（如我所想）是用更配得上一名信徒的猜测。让我们查看是否有什么来自古时的（原则）来支持我们。\par

在\BibleBook{}{出谷纪}中，梅瑟的那个姿势——藉祈祷与阿玛肋克作战，并坚持不懈直到\Quote{日落}——难道不是一个\Quote{晚站岗}吗？我们以为农的儿子若苏厄，在攻打阿摩黎人时，在他命令天象自身也\Quote{站岗}的那一天，他用了早餐吗？太阳\Quote{站住}在基贝红，月亮在阿雅隆谷；太阳和月亮\Quote{站住并停住，直到百姓向敌人报了仇}；太阳停在天空正中。此外，当天色渐晚，一日将尽时，以前和以后都没有这样的日子（当然，没有这样\Quote{长}的日子），\Quote{以致天主}，经上记着，\Quote{应允一个人的声音}——一个人，当然，与太阳同等，如此长久地坚持他的职守——一种甚至比\Quote{晚}更长的站岗。\par

无论如何，撒乌耳自己在交战时，明明地\Quote{命令}了这一职守：\Quote{谁若在晚上以前，在我向敌人报仇以前，吃了什么东西，是可咒骂的！}于是他的全体百姓都没有尝食物，而整片大地却在用早餐！天主还如此庄严地确认了那命令站岗的敕令，以至于撒乌耳的儿子约纳堂，虽然他在不知晓禁食延至很晚的情况下尝了一点蜂蜜，也立即被抽签定为有罪，并因百姓的祈祷才得以免除惩罚：因为他被定为贪食之罪，尽管是简单的一种。再者，达尼尔在达理阿王元年，身穿苦衣，禁食，向天主认罪时，说：\Quote{我还在祈祷发言的时候，在最初异象中我所见过的那人，迅速飞行，约在晚祭的时刻，来到我面前。}这将是一种\Quote{晚}站岗，它禁食\Quote{直到晚上}，并向上主献上一份更丰腴的祈祷祭！\par

% structure: p0038 source=h3 target=section reason=relative-heading-rank; base=h3

\section{第十一章 论对\Quote{人的权威}应有的尊重，以及论\Quote{异端}和\Quote{伪先知}的指控}

但我相信，所有这些（事例）要么不为那些对我们的行为感到不安的人所知，要么他们仅凭阅读而知道，却未加以仔细研究；这正是多数\Quote{无知的}自夸之徒——即\Quote{属血气的人}——的情况。这就是为什么我们径直穿越了禁食、干食、站岗的各种不同类别：为的是，当我们根据在旧约与新约中所发现的材料，叙述戒除、减少或推迟食物的虔诚遵守所带来的益处时，我们可以反驳那些把这些事贬低为空虚遵守的人；另一方面，当我们同样指出它们始终占据的宗教职责等级时，我们可以谴责那些把它们指责为标新立异的人：因为那一直存在的事并非新事，那有用的事也并非空虚。\par

然而，问题仍然摆在我们面前：有些规诫是天主为人类所命令的，已使这实践成为法律上具有约束力的；有些则是人为天主所奉献的，已履行了某种许愿的义务。然而，即使是许愿，一旦被天主接纳，也因接纳者的权威而成为未来的法律；因为凡赞同已行之事者，便已命令今后如此行。因此，从这一点来看，对立一方的争辩也被堵住了，他们声称：\Quote{若这是属神的声音制定了你们的这些庆节，那便是假预言；若这是人的臆测发明了它们，那便是异端。}因为他们一方面指责古代治理模式运行的形式，同时又从该形式中抽出论据来反击（我们），而古代治理模式的反对者自己也能反过来用这些论据回击，他们将不得不要么拒绝这些论据，要么承担这些被证实的义务（他们所抨击的）：必然如此；主要因为这些被抨击的义务，无论其制定者是谁，是属神的人还是普通信徒，都指向与古代治理模式相同的天主的荣耀。因为毫无疑问，在我们这些唯一创造主天主及其基督的司铎看来，异端和假预言都将因神性的差异而被判断：就此而言，我毫无偏袒地维护这一方，任凭对手选择任何立场来辩论。\Quote{那是魔鬼的灵，}你这样说，哦，属魂的人。而它怎能命令属于我们天主的义务，并命令只向我们的天主奉献呢？要么争辩说魔鬼与我们的天主合作，要么就让护慰者被视为撒殚。但你断言这是\Quote{人间的假基督}：因为若望正是以此名称呼异端者。而无论他是谁，他怎么能在我们基督的（名）下将这些义务指向我们的主呢？尽管假基督也曾（自称）走向天主，但却是反对我们的基督而出去。那么，你认为圣神在哪一边被证实存在于我们中间？是当他命令时，还是当他赞同时，我们的天主所经常命令和赞同的事？但你又一次为天主设立界线，如关于恩宠，同样关于纪律；如关于神恩，同样关于庆节：以至我们的实践被认为像祂的恩惠一样止息了；而你否认祂继续施加义务，因为在这种情况下，\Quote{法律及先知到若翰为止。}剩下的就是你彻底放逐祂，就\Emph{你}而言，祂是如此无所事事。\par

% structure: p0041 source=h3 target=section reason=relative-heading-rank; base=h3

\section{第十二章  反对属魂的人及其放纵之需要的某种抗议}

因为，此时，在这方面以及其他方面，\Quote{你们在财富和饱足中为王}——并不侵犯那些禁食所减少的罪，也不需要那些干食所强求的启示，也不理解那些站岗所驱散的你们自己的战争。即使从若翰时代起护慰者已沉默不语，我们自己也会成为自己的先知，主要为此缘故：我现在不是说以我们的祈祷平息天主的愤怒，或获得祂的保护或恩宠；而是通过预先防范来确保\Quote{末世}的道德立场；命令各种\OriginalTerm{谦卑之心}{\GreekText{ταπεινοφρόνησις}}，因为我们必须习惯于监狱，操练饥渴，以及获取忍受食物缺乏和为此焦虑的能力：为使基督徒进入监狱时，仿佛他刚从监狱出来一样——在那里承受的并非惩罚，而是纪律，并非世界的折磨，而是他自己惯常的实践；并以更大的信心离开拘留所面对（最终）战斗，因为他身上没有任何有罪的对肉体的虚假牵挂，从而使酷刑甚至没有材料可下手，因为他身穿干枯皮肤的铠甲，套上角质的护甲迎接利爪，他的血液已如灵魂的行李先被发送（上天），灵魂本身也加速（追随），通过频繁的禁食，已经获得了对死亡的极亲密认识！\par

显然，\Emph{你们的}习惯是在监狱中为不可靠的殉道者提供小厨房，生怕他们错过惯常的习俗，厌倦生命，（并）因新奇的禁欲纪律而跌倒；（一种纪律）连那位著名的普里斯提努斯——你们的殉道者，而不是\Emph{基督的}殉道者——也从未接触过：他——由于\Quote{自由拘留}（现今流行的）所提供的便利而长期被塞满，我想，他还对所有的浴室（仿佛它们比圣洗更好！）和所有纵欲的退隐处（仿佛它们比教会的更隐秘！）以及此世的所有诱惑（仿佛它们比永生的更有价值！）负有义务，因此不愿死——在审判的最后一天，正午时分，你们用搀药的酒作为解毒剂预先给他下药，使他如此衰弱，以至于当被几把利爪搔痒时——因为他的醉酒使这感觉像搔痒——他再也不能回答审讯官询问他\Quote{他承认谁为主}；于是他因这沉默被施以刑架，除了打嗝和嗳气什么也说不出，就在背教的行径中死去！这就是为何那些传讲清醒的人是\Quote{假先知}；这就是为何那些实践清醒的人是\Quote{异端者}！那么还犹豫什么，不信护慰者——你们在孟他努身上否认祂——存在于阿皮西乌斯身上呢？\par

% structure: p0044 source=h3 target=section reason=relative-heading-rank; base=h3

\section{第十三章  论属魂之人的矛盾}

你规定这信德有它的庆节，是圣经或祖先传统所\Quote{指定}的；并且为了革新的非法性，不得再添加任何遵守的附加。如果你能，就站在那个立场上。因为，看，我控告你在逾越节之外，在\BibleQuote{新郎被带走}的那些日子之外，也进行禁食；并且插入\OriginalTerm{站礼}{Station}的半禁食；而你（我发现）有时只靠面包和水生活，当每个人看来合适时（这样做）。简而言之，你回答说\Quote{这些事是出于选择，而不是出于命令而做的}。因此，你已改变了立场，通过超越传统，承担了那些没有被\Quote{指定}的遵守。但是，允许你自已的选择，却不允许天主的命令，这是什么行为？人的意志比天主的能力更有自由吗？我记得我是自由于\Emph{世界}，而非自由于天主。因此，我的职责是在没有外部提示的情况下，向我的主表达敬意，而祂命定的则是另一回事。我不但应当乐意服从祂，而且还应当讨祂喜悦；因为前者是我对祂命令的回应，后者是我自己选择的。\par

但对我来说，主教们也为全体教会团体发布禁食令，这已足够成为惯例；我指的不是为了收集施舍捐款的特殊目的，就像你们乞丐式的作风那样，而是有时也出于教会关怀的某些特定原因。因此，如果你们因一个人的命令而实践\OriginalTerm{谦卑}{\GreekText{ταπεινοφρόνησις}}，并且全体一致，那么为什么在我们的情况下，你们却指责我们禁食、\OriginalTerm{干食斋戒}{xerophagies}和\OriginalTerm{站礼}{Station}的合一呢？——除非，也许我们违反的是元老院的法令和皇帝反对\Quote{集会}的命令！圣神在祂选择的地方并通过祂选择的人宣讲时，常会预见到教会将面临的诱惑或世界将遭遇的灾祸，以\OriginalTerm{护慰者}{Paraclete}（即为通过祈祷赢得审判者而辩护的）的身份，颁布这类遵守的命令。例如，在当下，为了实践节制和斋戒的纪律：我们这些接受祂的人，也必须遵守祂当时所做的规定。看看犹太人的历法，你会发现，所有后代都世代谨慎地遵守给予祖先的诫命并非新奇之事。此外，在希腊各省，在特定地点举行从普世教会聚集的会议，通过这些会议，不仅所有更深层的问题为共同利益得到处理，而且整个基督徒名字的实际代表也受到极大的敬礼。（这是多么值得的事！在信德的护佑下，人们从四方聚集到基督那里！\BibleQuote{看，兄弟们同居共处，多么美好，多么快乐！}你们不太容易唱这首圣咏，除非你们与一大群人共进晚餐！）但那些会议首先通过站礼和禁食的操练，知道什么是\BibleQuote{与哭泣的人一同哭泣}，然后终于\BibleQuote{与欢乐的人一同欢乐}。如果我们在我们各自的省份，但在精神上彼此同在，遵守那些我们现在的论述所辩护的庆节本身，那就是圣事的法律。\par

% structure: p0047 source=h3 target=section reason=relative-heading-rank; base=h3

\section{第十四章。回复关于\Quote{效法迦拉达人}的指控。}

因此，为这些事遵守\BibleQuote{时节}，并遵守\BibleQuote{日子、月份、年份}，我们\Emph{效法迦拉达人}。显然，如果我们遵守\Emph{犹太人的}仪式、\Emph{法律的}庆节，我们就是效法迦拉达人。因为那些宗徒反对，他压制了已被埋葬在基督里的旧约的延续，并确立了新约的延续。但是，如果基督里有了新造，我们的庆节也必须是新的；否则，如果宗徒完全废除了所有\Quote{时节、日子、月份、年份}的虔诚，为什么我们在第一个月以年度循环庆祝逾越节？为什么在随后的五十天里我们充满喜乐地度过？为什么我们将一周的第四日和第六日奉献给站礼，将\OriginalTerm{预备日}{preparation-day}奉献给禁食？无论如何，你们有时甚至把站礼延长到安息日——除了在逾越节时期，根据其他地方给出的理由，安息日绝不可作为禁食日。在我们这里，无论如何，每一天也同样由普通的祝圣来庆祝。因此，在宗徒眼中，那区别原则——（如他所做的）区分\BibleQuote{新旧事物}——将不会是可笑的；但是（在此情况下）可笑的将是你们的不公，你们一边用\Emph{古代}的形式嘲弄我们，一边却用\Emph{新颖}的罪名控告我们。\par

% structure: p0049 source=h3 target=section reason=relative-heading-rank; base=h3

\section{第十五章。论宗徒关于食物的言论。}

宗徒同样斥责那些命令禁戒肉食的人，但他这样做是出于圣神的预见，预先谴责了那些将强制永远禁食以达致毁灭和轻视造物主之工的异端者；就如我在马西昂、塔提安或今日的毕达哥拉斯派异端者尤皮特身上所见的，而非在护慰者身上。因为我们的\Quote{禁戒肉食}是多么有限！一年中两星期的干食（且非全部——即安息日和主日除外）我们奉献给天主；禁戒那些我们不\Emph{拒绝}而只是\Emph{推迟}的事物。此外，当写信给罗马人时，宗徒现在给了\Emph{你们}一记直击，你们这些此规的诋毁者：\BibleQuote{不要为了食物，}他说，\BibleQuote{拆毁天主的工作。}什么\BibleQuote{工作？}就是关于他所说的：\BibleQuote{不吃肉、不喝酒是好的：} \BibleQuote{因为在这些事上服侍的人，是令我们的天主喜悦和悦纳的。} \BibleQuote{有人相信万物都可吃；但软弱的人只吃蔬菜。吃的人不可轻视不吃的人。你是谁，竟论断别人的仆人？} \BibleQuote{吃的人和不吃的人，都感谢天主。}但是，既然他禁止将\Emph{人的}选择作为争论的缘由，那么\Emph{神圣的}选择就更当如此了！因此他知道如何责备某些限制和禁戒食物的人，他们因轻视而非本分而禁食；但赞许那些为荣耀而非侮辱造物主而这样做的人。并且如果他\BibleQuote{将肉市的钥匙交给了你们}，允许吃\BibleQuote{万物}以确立对\BibleQuote{祭偶像之物}的例外，他仍未将天主的国纳入肉市：\BibleQuote{因为，}他说，\BibleQuote{天主的国不在于吃喝；}并且\BibleQuote{食物不能使我们受天主称许}——并非要你们认为这话是关于\Emph{干燥}的饮食，而是关于丰盛和精心烹调的；因为当他补充说：\BibleQuote{我们若吃了，也不多余；若不吃，也不缺少，}他的话的语气更适合（如其所是）你们而非我们，你们认为若吃了就\BibleQuote{多余}，若不吃就缺少，并因此贬低这些实践。\par

你们将主随意\BibleQuote{又吃又喝}的事实解释为偏袒自己的私欲，这多么不相称啊！但我想祂也必定\BibleQuote{禁食过}，因为祂曾宣告：不是\BibleQuote{饱足的}，而是\BibleQuote{饥饿和口渴的人是有福的}；祂惯于宣称\BibleQuote{食物}不是门徒所想的，而是\BibleQuote{彻底完成圣父的工作}；教导\BibleQuote{要为那存留到永生的食物劳碌}；同样在通常的祈祷中命令我们求\BibleQuote{面包}，而非亚他罗的财富。如此，依撒意亚也\Emph{没有}否认天主\BibleQuote{拣选了}\Emph{一种}\BibleQuote{禁食}；而是详细说明了祂\Emph{没有}拣选的禁食\Emph{种类}：\BibleQuote{因为在你们禁食的日子，}他说，\BibleQuote{你们仍寻自己的意愿，并且你们暗中刺痛一切服从你们的人；否则你们禁食是为了虐待和争吵，并用拳头打人。不是\Emph{这样的}禁食我选择了；}而是祂随后附加的那种，并且藉着附加没有废除，反而确认了。\par

% structure: p0052 source=h3 target=section reason=relative-heading-rank; base=h3

\section{第十六章 圣经中天主对放纵者的审判实例；并诉诸外邦人的实践}

因为即使祂\Emph{更喜爱}\Quote{正义的作为}，但仍不是没有祭献，即一个以禁食所压伏的灵魂。无论如何，祂是那位既不喜悦放纵食欲的百姓，也不喜悦放纵的司祭或先知的天主。直到今日，\BibleQuote{贪欲的坟墓}仍存留着，那里贪求\BibleQuote{肉}的百姓，因狼吞虎咽鹌鹑而不消化，患了霍乱，都被埋葬。厄里在圣殿门前折断颈项，他的儿子们阵亡，他的儿媳死于分娩：因为无羞耻之家骗取肉祭，确应受天主这样的打击。舍玛雅，一位\BibleQuote{天主的人}，在预言了雅洛贝罕王引入偶像崇拜的结果之后——在王的手枯干又立即复原之后——在祭台裂为两半之后——因这些征兆被王邀请回家作为酬报，明确拒绝了（因天主曾禁止他）在当地碰任何食物；但不久之后从另一位谎称自己是先知的老者那里鲁莽地取了食物，便按照当时在桌前发出的天主的话，被剥夺了葬于祖坟的权利。因为他在路上被一头冲来的狮子扑倒，葬在陌生人中；如此为他破禁食的惩罚付出了代价。\par

这些事将成为对民众和主教，甚至属灵之人的警告，以防他们可能曾犯过口腹之欲的放纵。不，即使在阴间，劝诫也未停止说话；我们在那富有的宴乐者身上看到宴饮受折磨；在穷人身上看到斋戒得安舒——（正如宴饮和斋戒一样）——有\BibleQuote{梅瑟和先知}作为教师。岳厄尔也曾大声疾呼：\BibleQuote{祝圣一个斋戒和一个宗教仪式}；他那时已预见到其他宗徒和先知会认可斋戒，并会宣讲对天主的特别侍奉。因此，连那些以装扮、在圣所中装饰并向\Emph{偶像}在每个时辰敬礼来追求\Emph{偶像}的人，也被说成是在\Emph{侍奉}它们。但不仅如此，外邦人也认同各种形式的\OriginalTerm{谦逊}{\GreekText{ταπεινοφρ}\'{\GreekText{ο}}\GreekText{νησις}}。当天空僵硬、年景干旱时，公开宣告命令赤足游行；官员们脱下紫袍，翻转权杖，祈祷，献祭。此外，有些殖民地在这些特殊仪式之外，还以年度礼仪，身穿麻衣、撒灰，向他们的偶像呈现恳求的急迫，而浴场和店铺关闭直到第九时辰。他们公共场合只有一处火——在祭坛上；连盘子里也没有水。我相信有一种尼尼微式的暂停营业！无论如何，犹太人的斋戒普遍举行；他们无视圣殿，在整个岸边、每个露天场所，长久地向天祈祷。而且，尽管他们以哀悼的服饰和装饰羞辱了职责，但他们仍实际表示了对禁食的信德，并叹息盼望迟迟未到的晚星来临以允许进食。但对我来说，你们通过将亵渎堆在\Emph{我们的}干食上，将其与伊西斯和库柏勒的贞洁并列，这已足够。我承认这种比较作为证据。因此，我们的干食将被证明是神圣的，因为魔鬼，神圣事物的模仿者，模仿它。谎言是建立在真理之上的；迷信是由宗教结合而成的。因此，\Emph{你们}更不虔诚，因为外邦人更顺从。简言之，他向偶像之神牺牲他的食欲；而\Emph{你们}却不愿向（真）天主如此行。因为对你们来说，你们的肚腹是神，你们的肺是圣殿，你们的胃是祭坛，你们的厨师是司铎，你们的香味是圣神，你们的调味品是属灵的恩赐，你们的打嗝是预言。\par

% structure: p0055 source=h3 target=section reason=relative-heading-rank; base=h3

\section{第十七章　结论}

如果我们要说实话，你们这些如此纵容食欲的人，你们是\Quote{老}的，你们正当地夸耀你们的\Quote{优先}；我总能认出厄撒乌——猎取野兽者的气味：你们如此无限制地热衷于捕捉田鸫，你们从最松懈纪律的\Quote{田野}而来，你们在精神上如此软弱。如果我向你们提供一份便宜的扁豆，用浓缩的葡萄汁染红，你们会立即卖掉你们所有的\Quote{首领地位}；在你们那里，\Quote{爱}在炖锅中显示其热情，\Quote{信德}在厨房中温暖，\Quote{望德}在侍者中锚定；但\Quote{爱}更加重要，因为它是你们年轻人与姐妹同寝的手段！正如我们都知道，食欲的附属品是好色和纵欲。宗徒也知道这种联合；因此，在预先说过\BibleQuote{不在醉酒和宴乐中}之后，他附加了\BibleQuote{也不在床笫和情欲中}。\par

对你们食欲的控告包括这一点：\Quote{双倍荣耀}被你们分配给主礼的长老（以双份的肉和饮料）；而宗徒给予他们\Quote{双倍荣耀}，既是\Emph{弟兄}又是\Emph{职员}。你们中谁在圣德上更高超，除非是那个更频繁赴宴、更奢华备办食物、更精通酒杯的人？你们只是灵魂和肉体的人，你们正当地拒绝属灵之事。如果先知们取悦于\Emph{这样的}人，那我的先知们就不是。那么，你们为什么不经常宣讲：\BibleQuote{我们吃喝吧，因为明天我们要死了？}正如\Emph{我们}毫不畏惧地勇敢命令：\Quote{我们斋戒吧，弟兄姐妹们，以免明天我们或许死去。}让我们公开维护我们的纪律。我们确知：\BibleQuote{属于肉性的人不能中悦天主}；当然，不是指处于肉身的\Emph{实体}中的人，而是处于肉身的\Emph{挂虑}、\Emph{情感}、\Emph{工作}、\Emph{意愿}中的人。消瘦不使我们不悦；因为天主赐予肉身不是按重量，正如祂赐予\BibleQuote{圣神不是按尺度}。或许，通过救恩的\BibleQuote{窄门}，更瘦的肉身更容易进入；更轻的肉身复活更快；在坟墓中更干的肉身体持更久。让奥林匹克拳击手和斗士们把自己塞饱。对身体的雄心适合于那些需要身体力量的人；然而他们也通过干食来加强自己。但我们的肌肉和筋骨是另一种，正如我们的竞赛也是另一种；我们\BibleQuote{的摔跤不是对抗血肉，而是对抗世界的权力，对抗邪恶的属灵势力}。对抗这些，我们应当以信德和精神的强壮，而非血肉的强壮来作对抗的站立。相反，一个过度喂养的基督徒对熊和狮子可能比天主更需要；只不过，即使面对野兽，他也应当练习消瘦。\par

\SourceRecord{Translated by S. Thelwall.；Ante-Nicene Fathers；Vol. 4.；Edited by Alexander Roberts, James Donaldson, and A. Cleveland Coxe.；Buffalo, NY: Christian Literature Publishing Co.,；1885.；<http://www.newadvent.org/fathers/0408.htm>.}
